% --------------------------------- Preamble -----------------------------------
\documentclass[11pt, letterpaper, twocolumn]{article}

\usepackage[T1]{fontenc} % use 8-bit fonts
\usepackage[margin=0.75in, columnsep=0.2in, % 3.4 in columns
    headsep=0.28in, footskip=0.42in]{geometry} % paper sizes
\usepackage{amsmath, amsfonts, amssymb} % improved math
\usepackage{bm} % bold math (must be after amsmath)
\usepackage{upquote} % upright single quotes in verbatim
\usepackage{pgfplots} % includes tikz, xcolor, graphicx
\usepackage[colorlinks=true, allcolors=gray]{hyperref}

\raggedbottom % do not force stretching text to fill the page
\setlength{\emergencystretch}{\hsize} % avoid overfull lines
\overfullrule=1mm % black band to show overfull lines
\pdfsuppresswarningpagegroup=1 % remove pdf image warnings
\frenchspacing % consistent spacing between words
\newcommand{\ul}[1] % vectors
    {{}\mkern1mu\underline{\mkern-1mu#1\mkern-1mu}\mkern1mu}
\DeclareSymbolFont{euler}{U}{eur}{m}{n} % symbol font
\DeclareMathSymbol{\PI}{\mathord}{euler}{25} % up pi
\pgfplotsset{compat=newest} % use newest pgfplot version
\newcommand{\EE}[2]{\ensuremath{{{#1}\!\times\!10^{#2}}}}
% ------------------------------------------------------------------------------

\title{2D Bayesian Filtering with Multimodal Dynamics}
\author{}
\date{}

\begin{document}
\maketitle

\section*{Description}

The overall goal of this project is to implement a Bayesian filter over a 2D
continuous space, estimating the location of a car with a measurement model and
a dynamic model. The dynamic model has no deterministic propagation, only a
dynamic noise PDF.

The state of the system is the two-dimensional position of the vehicle,
\[
x = \begin{bmatrix} x \\ y \end{bmatrix}
\]

Each time step, the car moves approximately 5~m. It can either move forward 5~m,
or mainly forward but to the left or right as well (forward 4~m, left or right
3~m). More specifically, the probability of moving left is 0.3, right is 0.3,
and straight is 0.4. Around each movement, there is an isotropic (same in all
directions) Gaussian with variance of 0.1~m$^2$. These possible motions are shown here: \\
\includegraphics[width=\linewidth]{../figures/fig_motions.png}

The measurement also will be an isotropic Gaussian distribution with variance
4~m$^2$.

The initial location of the car is a Gaussian random variable, centered at (0,
0), with an isotropic covariance matrix of 0.6~m$^2$.

\section*{Discretization and Representation}

The continuous state space is approximated using a discrete grid with resolution
$dx = 0.2$ m. The domain is chosen to encompass both the true trajectory and
measurement data, with additional margin to account for uncertainty.

A mesh grid $(X,Y)$ is constructed to represent all possible positions.
Probability density functions are evaluated over this grid.

The initial state, process noise, and measurement likelihood are all
modeled as Gaussian distributions evaluated on the grid using a 
two-dimensional normal distribution.

\section*{Bayesian Filtering}

The state distribution is propagated using a recursive Bayesian filtering framework.

\textbf{Measurement Update:} The prior distribution is updated using the measurement
likelihood:
\[
p(x) \propto p(z{\mid}x) p(x)
\]

\textbf{Time Propagation:} The updated distribution is propagated using convolution
with the process noise distribution:
\[
p(x_{k+1}) = p(x_k)*p(\nu)
\]

The convolution is implemented numerically using a two-dimensional discrete convolution,
which approximates the continuous propagation of uncertainty.

\section*{Multimodal Dynamics}

The motion model consists of three possible movements: left, straight, and right, each
with associated probabilities. The overall process noise distribution is constructed as
a weighted sum of Gaussian distributions centered at the corresponding motion offsets.

This results in a multimodal distribution that captures multiple possible future states,
which is then used in the convolution step for time propagation.

\section*{Results}

\subsection*{Dynamics Only}

Figure~\ref{fig:dynamics} shows the evolution of the prior distribution when
only the dynamics model is applied. The distribution spreads and shifts over
time according to the multimodal motion model.

\begin{figure}[h]
    \centering
    \includegraphics[width=0.8\linewidth]{../figures/fig_dynamics_only.pdf}
    \caption{Evolution of the state distribution under the dynamics model.}
    \label{fig:dynamics}
\end{figure}

\pagebreak

\subsection*{Full Estimation}

Figure~\ref{fig:estimation} shows the full Bayesian filtering process, including
both prediction and measurement updates.

\begin{figure}[h]
    \centering
    \includegraphics[width=\linewidth]{../figures/fig_full_estimation.pdf}
    \caption{State estimation with measurement updates over time.}
    \label{fig:estimation}
\end{figure}

\subsection*{Observations}

We notice that the prior is often multimodal, but the update typcially has
a single peak. The prior is often multimodal because it represents the hypotheses 
about the objects movement, where it can either move left, straight, or right.
The probabilities are each of a similar magnitude, so the prior gives use
3 distinct peaks. When a new measurement arrives, it creates a single 
"spotlight" of evidence, which is centered around the new measurement.
During the update, the evidence multiplies the prior, amplifying the one
hypothesis it agrees with the most, suppressing the others. This is why
the posterior distribution usually collapses into a single peak.

\section*{Discussion}

The results illustrate how uncertainty evolves under both the dynamics and
measurement models. The multimodal motion model leads to a distribution
that can split into multiple regions, reflecting different possible
trajectories.

Measurement updates help concentrate the distribution around the true state,
reducing uncertainty over time. This demonstrates the balance between process
uncertainty and measurement information in Bayesian estimation.

The use of convolution provides a natural way to propagate uncertainty in
continuous spaces, forming a foundation for more advanced estimation techniques.

\end{document}